\label{chap:p3-83-46-godel_turing}

\section{Décidabilité effective de chaque fenêtre finie de multiplicités projectives}

\begin{proposition}
Une profondeur \(N\) étant fixée, si les domaines sont finis et les observables
atomiques sont calculables, toute formule du langage fini
\(L_N\) est sémantiquement décidable.
\end{proposition}

\begin{proof}
Les formules atomiques se calculent ; les connecteurs se décident par des tables de
vérité ; chaque quantificateur parcourt un domaine fini. On conclut par
induction sur la complexité syntaxique.
\end{proof}

La décidabilité de chaque fenêtre ne fournit pas nécessairement un seul
algorithme qui décide une propriété de la tour complète.

\section{Tour effective associée au calcul de multiplicités projectives}

Soit \(T_e\) la machine de Turing d'indice \(e\). Pour chaque \(n\), nous définissons
\[
 X_n^e
 =
 \{b_n\}
 \cup
 \{c_n:T_e\text{ ne s'est pas arrêté dans ses premiers }n\text{ pas}\}.
\]
Les liens conservent les étiquettes existantes :
\[
  p_{n+1,n}(b_{n+1})=b_n,
  \qquad
  p_{n+1,n}(c_{n+1})=c_n.
\]
La tour visible possède un seul symbole \(v_n\), et
\[
  q_n(b_n)=q_n(c_n)=v_n.
\]

\begin{theorem}[Multiplicité globale et arrêt]
\label{p3-83-c46-s1:thm:halt-fibre}
\[
 \#q_\infty^{-1}(\boldsymbol v)
 =
 \begin{cases}
 1,&T_e\text{ s'arrête},\\
 2,&T_e\text{ ne s'arrête jamais}.
 \end{cases}
\]
Par conséquent, il n'existe pas d'algorithme total qui décide l'unicité de la fibre
globale pour toutes les tours effectives de cette famille.
\end{theorem}

\begin{proof}
La branche \(\boldsymbol b=(b_n)\) existe toujours. Si \(T_e\) s'arrête au
pas \(h\), les états \(c_n\) disparaissent à partir de \(h\) et ne forment pas une
section globale. Si elle ne s'arrête jamais, \(\boldsymbol c=(c_n)\) est une deuxième
section. Un décideur de l'unicité résoudrait \(\HALT\), en contradiction avec
\citep{turing1936}.
\end{proof}

On définit
\[
  U_e:\#q_\infty^{-1}(\boldsymbol v)=1.
\]
Alors \(U_e\iff T_e\downarrow\).

\begin{theorem}[Incomplétude uniforme]
\label{p3-83-c46-s1:thm:uniform-incompleteness}
Soit \(\mathsf T\) une théorie récursivement axiomatisable qui formalise les
tours \(X^e\) et est sémantiquement correcte pour \(U_e\) et
\(\neg U_e\). Alors \(\mathsf T\) ne décide pas tous les \(U_e\).
\end{theorem}

\begin{proof}
Si elle décidait chaque \(U_e\), pour un indice \(e\) on énumérerait en parallèle
les preuves de \(U_e\) et \(\neg U_e\) jusqu'à en trouver une. Par correction,
la première apparaîtrait exactement lorsque la machine s'arrête et la seconde
lorsqu'elle ne s'arrête pas. On obtiendrait un décideur de \(\HALT\).
\end{proof}

\begin{figure}[H]
\centering
\begin{tikzpicture}[>=Latex,node distance=11mm and 15mm,
  box/.style={rounded corners,draw=hmtblue,fill=hmtlightblue,
    minimum width=39mm,minimum height=13mm,align=center},
  warn/.style={rounded corners,draw=hmtred,fill=hmtlightred,
    minimum width=39mm,minimum height=13mm,align=center}]
  \node[box] (sim) {simulation finie de \(T_e\)};
  \node[box,below=of sim] (levels) {\(X_n^e\) calculable\\pour chaque \(n\)};
  \node[box,below=of levels] (fibre)
    {\(\#q_\infty^{-1}(v)=1\iff T_e\downarrow\)};
  \node[warn,below left=of fibre] (undec) {sans décideur uniforme};
  \node[warn,below right=of fibre] (inc)
    {théorie effective correcte\\incomplète pour les \(U_e\)};
  \draw[->,thick,hmtblue] (sim)--(levels);
  \draw[->,thick,hmtblue] (levels)--(fibre);
  \draw[->,thick,hmtred] (fibre)--(undec);
  \draw[->,thick,hmtred] (fibre)--(inc);
\end{tikzpicture}
\caption{De la décision finie à l'incomplétude uniforme.}
\label{p3-83-c46-s1:fig:godel-turing}
\end{figure}

\section{Complétude finie des niveaux HMT et domaine exact de l'obstruction de Gödel--Turing}

Le nom académique utilisé est :
\[
\boxed{\text{Théorème HMT d'incomplétude uniforme pour les multiplicités projectives}.}
\]
On emploie également la dénomination abrégée « troisième résultat de Gödel dans
HMT ». La preuve
est de Gödel--Turing : elle ne reproduit pas la diagonale autoréférente de
\cite{goedel1931}, mais réduit HALT à la multiplicité d'une fibre.

\begin{remark}[Conditions de validité]
Le théorème ne prouve pas l'incohérence de ZFC. Il n'empêche pas non plus de démontrer
directement qu'une tour particulière possède une section unique. Il empêche une
procédure universelle d'unicité pour toutes les tours effectives.
\end{remark}

% REMATE_LOCAL_GODEL_CONCLUSION
\Needspace{25\baselineskip}
\section{Relation exacte avec les \(9\) phases du Topological Phase Kernel}

Les \(9\) phases du TPK apportent un algorithme concret qui, avec ses lecteurs déclarés,
partitionne le bloc de raffinement en \(729\) sous-cylindres et sélectionne l'un d'eux
à chaque étape finie.
Le passage à la section infinie utilise un théorème global de non-vacuité,
de compatibilité et d'unitarité.

Le résultat de Gödel--Turing dit :
\[
\text{il n'existe pas de décideur d'unicité pour toute tour effective}.
\]
Il n'affirme pas :
\[
\text{aucune tour concrète ne peut avoir de preuve uniforme propre}.
\]
Et une preuve particulière ne fournit pas le décideur universel interdit.


\begin{remark}[Conclusion structurelle]
Décision finie, existence de la limite, unicité d'une instance et décision
uniforme de toutes les instances sont quatre niveaux. La clôture holographique
exige de les déclarer au lieu de transférer automatiquement une propriété d'un
niveau à un autre.
\end{remark}

\begin{remark}[Séparation des résultats et portée]
La tour réductrice, l'équivalence \(U_e\iff T_e\downarrow\) et le théorème
métamathématique sont des \emph{résultats établis dans les sections précédentes}. La position centrale au sein
de la clôture de l'infini est une organisation de l'exposé qui ne modifie ni le contenu logique ni le domaine du théorème.
\end{remark}
